% =====================================================================
% CHAPTER 6: Q2 - DO GENERATED REPLIES LOOK LIKE REAL EMAIL? budget: 5 pages
% ROLE: Answer Q2. Story of the chapter:
%   models stay on topic -> but are easy to tell from real email, by style ->
%   judges still rate them as plausible -> and they are less diverse.
% =====================================================================
\chapter{Q2: Do Generated Replies Look Like Real Email?}
\label{ch:q2}

% P0 KEY: Q2 uses the paired Design A: real and generated replies to the
%         same 183 emails.
% SUPPORT:
%   - Exceptions: rubric (judge-swap set) and Self-BLEU within a scenario (Design B).

\section{Mirroring}
% budget: 0.5 pages

% P1 KEY: Copying the sender's words is mainly a weakness of the small model.
% SUPPORT:
%   - Borrowed words (same length): Llama 0.565 (19.1% flagged), 20b 0.373,
%     120b 0.435, qwen 0.298, real 0.294.
%   - All models borrow at least as much as real writers.
%   - FIG (optional): nvidia_four_models_mirroring.png. SRC: RO 3.1.

\section{Telling real from generated}
% budget: 2 pages

% P2 KEY: A simple word classifier tells real from generated replies almost
%         always.
% SUPPORT:
%   - AUC, same length: Llama 0.881, 20b 0.970, 120b 0.962, qwen 0.982.
%   - Length alone explains it only for Llama (0.837).
%   - SRC: RO 3.2.

% P3 KEY: The models are given away by stock phrases.
% SUPPORT:
%   - "I'll" in 59-61% of gpt-oss replies vs. 5-6% of real ones; "let"
%     (Let me know); "review"; "team".
%   - Real replies: sign-off names, "attached", "office".
%   - A larger model was not harder to spot.

% P4 KEY: In meaning space the replies stay tied to the email they answer.
% SUPPORT:
%   - nomic embeddings: AUC 0.82-0.94 (same length).
%   - "Tracks its own email" 0.81-0.86.
%   - SRC: RO 3.3.

% P5 KEY: Style alone separates them even better. The gap is how replies are
%         written, not what they say.
% SUPPORT:
%   - StyleDistance: AUC 0.97-0.99; "tracks its own email" falls to about 0.5,
%     as intended.
%   - Llama moves most (0.82 -> 0.97).
%   - FIG: embedding_style_vs_nomic_auc.png.

\section{Rubric scores from judges}
% budget: 1 page

% P6 KEY: Judges rate the larger models' replies as high quality and plausible.
% SUPPORT:
%   - Overall mean 4.4-4.6 of 5 under every judge.
%   - Llama depends on the judge (4.28 own judge; 3.48-3.60 gpt-oss judges),
%     mainly on role_consistency and contextual_fit.
%   - TAB: overall mean by writer x judge. Per-item table in the appendix.
%   - SRC: RO 3.4.

% P7 KEY: This disagrees with the classifier: "reads as plausible" is not the
%         same as "cannot be told apart from real".
% SUPPORT:
%   - Same replies: rubric ~4.5, classifier AUC ~0.97.
%   - The judge saw only the reply, not the role or message (Ch. 7).

\section{Diversity}
% budget: 1 page

% P8 KEY: All models are less diverse than real email.
% SUPPORT:
%   - Self-BLEU generated 0.51-0.57 vs. real 0.27-0.31 (same length), all p<.001.
%   - Fits P3: repeated stock phrases.
%   - FIG: self_bleu_vs_real.png. SRC: RO 3.5.

% P9 KEY: A different persona changes qwen's wording a lot and gpt-oss-120b's
%         very little.
% SUPPORT:
%   - Within one scenario: qwen 0.17, 20b 0.30, Llama 0.38, 120b 0.60.
%   - Model size does not decide diversity.

\section{Summary}
% budget: 0.5 pages

% P10 KEY: The answer to Q2, and the bridge to Q3.
% SUPPORT:
%   - On topic, but easy to spot by style, and less diverse.
%   - Judges call them plausible. So can the judge be trusted? -> Ch. 7.
